\section{Factorización de la Covarianza}

\begin{frame}{Factorización RS de la Covarianza}
  \begin{columns}
    \begin{column}{0.5\textwidth}
      Para garantizar $\boldsymbol{\Sigma}$ sea SPSD y optimizable:
      \[
        \boldsymbol{\Sigma} = \mathbf{R}\mathbf{S}\mathbf{S}^\top\mathbf{R}^\top
      \]
      \begin{itemize}
        \item $\mathbf{R} \in SO(3)$: matriz de rotación
        \item $\mathbf{S} = \mathrm{diag}(s_1, s_2, s_3)$, $s_k > 0$: escala
      \end{itemize}
      \vspace{0.5em}
      \textbf{Garantía SPSD:}
      \[
        \mathbf{v}^\top\boldsymbol{\Sigma}\mathbf{v} = \|\mathbf{S}^\top\mathbf{R}^\top\mathbf{v}\|^2 \geq 0
      \]
    \end{column}
    \begin{column}{0.5\textwidth}
      \begin{exampleblock}{Ventajas de la factorización}
        \begin{itemize}
          \item Optimización sin restricciones sobre $\mathbf{R}$ y $\mathbf{S}$
          \item $\boldsymbol{\Sigma}$ siempre válida por construcción
          \item Interpretación geométrica clara: rotación + escala de un elipsoide
        \end{itemize}
      \end{exampleblock}
    \end{column}
  \end{columns}
\end{frame}

\begin{frame}{Cuaterniones para la Rotación}
  \begin{columns}
    \begin{column}{0.5\textwidth}
      La rotación se almacena como cuaternión unitario:
      \[
        \mathbf{q} = (q_w, q_x, q_y, q_z), \quad \|\mathbf{q}\| = 1
      \]
      \textbf{Ventajas sobre ángulos de Euler:}
      \begin{itemize}
        \item Solo 4 parámetros (vs 9 de la matriz)
        \item Evita el \emph{gimbal lock}
        \item Interpolación suave (SLERP)
      \end{itemize}
    \end{column}
    \begin{column}{0.5\textwidth}
      \textbf{Matriz de rotación inducida:}
      \[
        \mathbf{R} = \begin{pmatrix}
          1-2(q_y^2+q_z^2) & 2(q_xq_y-q_wq_z) & 2(q_xq_z+q_wq_y)\\
          2(q_xq_y+q_wq_z) & 1-2(q_x^2+q_z^2) & 2(q_yq_z-q_wq_x)\\
          2(q_xq_z-q_wq_y) & 2(q_yq_z+q_wq_x) & 1-2(q_x^2+q_y^2)
        \end{pmatrix}
      \]
      \vspace{0.3em}
      \textit{Renormalizado tras cada paso de gradiente.}
    \end{column}
  \end{columns}
\end{frame}
